\subsection{Ungeeignete und eingeschränkt geeignete Projekttypen}
\label{subsec:ungeeignete-projekttypen}

Eingeschränkt geeignet sind Vorhaben, deren Zugriffskanal zwar einem Profil
entspricht, deren Kern aber deutlich andere technische Eigenschaften
verlangt. Bei umfangreicher Betriebssystemintegration, plattformspezifischer
UI, systemnahen Diensten oder Treibern reicht die vorhandene Tauri-Hülle nicht
als Nachweis der Eignung; native Kommandos, Berechtigungen und Tests müssten in
erheblichem Umfang produktspezifisch ergänzt werden. Gleiches gilt für eine
zwingend abweichende Infrastruktur, etwa eine stark verteilte
Microservice-Landschaft, spezielle Event-Streaming- oder HPC-Anforderungen.
Hier ist zu prüfen, ob die Baseline als abgegrenzter Client oder Dienst
weiterverwendet werden kann oder eine andere Architektur zweckmäßiger ist.

Nicht zum primären Zielbereich gehören harte Echtzeitsysteme, deren Korrektheit
von garantierten Antwortzeiten abhängt, grafikintensive Spiele mit Bedarf an
einer Game Engine sowie Firmware-, Mikrocontroller- und treiberzentrierte
Embedded-Systeme. Gewöhnliche Live-Aktualisierungen über HTTP oder WebSockets
sind davon zu unterscheiden: Sie machen eine Business-Anwendung noch nicht zu
einem Hard-Real-Time-System. Native iOS- und Android-Anwendungen sind ebenfalls
kein vorhandenes Standardprofil. Diese Abgrenzungen bedeuten nicht, dass jede
spielähnliche, mobile angebundene oder hardwarenahe Teilfunktion technisch
unmöglich wäre; ihr jeweiliger Kern wird durch die untersuchte Baseline jedoch
nicht bereitgestellt.

Hohe Sicherheits- oder Compliance-Anforderungen führen nicht automatisch zu
einer negativen Einstufung. Die Baseline allein weist aber weder
produktspezifische Schutzbedarfe noch Auditierung, Zertifizierung oder
Rechtskonformität nach. Der in
Kapitel~\ref{sec:qualitaetssicherung-verifikation} dokumentierte Stand belegt
zwar Compose-Validierung und Master-Image-Builds, nicht jedoch einen
Compose-Start oder produktive Bereitstellung
\parencite{kleiveist2026ci-core-run}. Linux- und macOS-Pakete wurden
unsigniert gebaut; der Windows-Paketpfad ist fehlgeschlagen
\parencite{kleiveist2026ci-desktop-run}. Signing und Notarisierung bleiben
produktspezifische Prüfgegenstände \parencite{kleiveist2026release}. Solche
Projekte erfordern daher zusätzliche Nachweise. Bei der Einsatzentscheidung
ist ferner zu berücksichtigen, ob im Team die für die gewählte Variante
benötigten TypeScript-, Python-, Rust-, Datenbank- und Deploymentkompetenzen
vorhanden sind.
